% Chapter 8, Section 1

\section{Gradient Descent Variants}
\label{sec:gd-variants}

\subsection{Intuition: Following the Steepest Downhill Direction}

Imagine standing on a foggy hillside trying to reach the valley. You feel the slope under your feet and take a step downhill. \textbf{Batch gradient descent} measures the average slope using the whole landscape (dataset) before each step: accurate but slow to gauge. \textbf{Stochastic gradient descent (SGD)} feels the slope at a single point: fast, but noisy. \textbf{Mini-batch} is a compromise: feel the slope at a handful of nearby points to get a reliable yet efficient direction. This trade-off underlies most practical training regimes.\index{optimization!gradient descent}\index{mini-batch}\gls{mini-batch}

Historical note: Early neural network training widely used batch gradient descent, but the rise of large datasets and GPUs made mini-batch SGD the de facto standard \cite{GoodfellowEtAl2016,Prince2023}.

\subsection{Batch Gradient Descent}

Computes gradient using entire training set:
\begin{equation}
\vect{\theta}_{t+1} = \vect{\theta}_t - \alpha \nabla_{\vect{\theta}} \frac{1}{n} \sum_{i=1}^{n} L(\vect{\theta}, \vect{x}^{(i)}, y^{(i)})
\end{equation}

Characteristics:\index{gradient descent!batch}
\begin{itemize}
    \item Guaranteed convergence to global minimum (convex) or local minimum (non-convex)
    \item Very slow for large datasets
    \item Deterministic updates
\end{itemize}

\subsection{Stochastic Gradient Descent (SGD)}\index{SGD@stochastic gradient descent}

Uses a single random example per update:
\begin{equation}
\vect{\theta}_{t+1} = \vect{\theta}_t - \alpha \nabla_{\vect{\theta}} L(\vect{\theta}, \vect{x}^{(i)}, y^{(i)})
\end{equation}

Characteristics:
\begin{itemize}
    \item Much faster per iteration
    \item Noisy updates can escape shallow local minima
    \item High variance in updates
    \item May not converge to exact minimum
\end{itemize}

\subsection{Mini-Batch Gradient Descent}\index{gradient descent!mini-batch}

Balances batch and stochastic approaches:
\begin{equation}
\vect{\theta}_{t+1} = \vect{\theta}_t - \alpha \nabla_{\vect{\theta}} \frac{1}{|\mathcal{B}|} \sum_{i \in \mathcal{B}} L(\vect{\theta}, \vect{x}^{(i)}, y^{(i)})
\end{equation}

where $\mathcal{B}$ is a mini-batch (typically 32-256 examples).

Benefits:
\begin{itemize}
    \item Reduces gradient variance
    \item Efficient GPU utilization
    \item Good balance of speed and stability
\end{itemize}

